
\chapter{Espacios compactos}
\section{Subespacios compactos y propiedades fundamentales}
Sea $X$ un espacio topol\'ogico
\begin{defi}
    Un \textbf{recubrimiento} de $X$ es una familia de abiertos $\{U_i\}_{i \in I}$ tal que $\displaystyle X = \bigcup_{i \in I} U_i$. 
\end{defi}
\begin{defi}
    Un espacio topol\'ogico $X$ es \textbf{compacto} si \underline{para todo} recubrimiento $\{U_i\}_{i \in I}$ existe un subrrecubrimiento finito, es decir, existen $i_1, \ldots, i_r \in I$ tal que $\{U_{i_1}, \ldots, U_{i_r}\}$ recubren $X$
\end{defi}
\begin{ejem}
$\bb R$ no es compacto. $\{U_n = (-n,n): n \in \bb N\}$ recubrimiento de $\bb R$. Si existen $n_1, \ldots, n_r \in \bb N$ tal que 
$$\bigcup_{j=1}^r U_{n_j} = (-n_1, n_1)\cup \cdots \cup (-n_r,n_r) = (-n,n) \neq \bb R$$
con $n = \max\{n_{i_j}\}$ entonces no recubren.

Otro ejemplo es el recubrimiento $\left\{V_n = \left(n- \frac{1}{|n|}, (n+1) + \frac{1}{|n|} \right)\right\}_{n\in \bb Z}$.
\end{ejem}
\begin{defi}
    Un \textbf{subespacio} topol\'ogico $Y \subset X$ es \textbf{compacto} si lo es como espacio topol\'ogico, es decir, si $(Y, \mc T_X|_Y)$ es compacto.
\end{defi}
\begin{prop}
$Y \subset X$ es subespacio compacto $\iff$ para toda familia de abiertos $\{U_i\}_{i \in I}$ de $X$ tal que $Y \subset \bigcup_{i \in I} U_i$ existen $\{i_1, \ldots, i_r\}\subset I$ tal que $Y \subset U_{i_r} \cup \cdots U_{i_r}$.
\end{prop}
\begin{dem} \,
\begin{itemize}
    \item[$\Rightarrow)$] $Y\subset \displaystyle \bigcup_{i \in I} U_i$, entonces $Y= \displaystyle \bigcup_{i \in I} (U_i\cap Y)$ y $ U_i\cap Y \in \mc T_X|_Y$, luego por ser $Y$ compacto, y tener un subrrecubrimiento finito, entonces
    $$Y = \bigcup_{j=1}^r U_{i_j} \cap Y \subset \bigcup_{j=1}^r U_{i_j}$$
    Por tanto, $\{U_{i_1}, \ldots, U_{i_r}\}$ recubre a $Y$.
    \item[$\Leftarrow)$] Sea $\{ V_i \}_{i\in I}$ recubrimiento de $Y$ por abiertos de $Y$, entonces existen $\{U_i\}_{i\in I}$ abiertos de $X$ tal que $U_i \cap Y = V_i\, \forall i \in I$. 
    $$Y = \bigcup_{i\in I} V_i = \bigcup_{i\in I} (U_i \cap Y) \subset \bigcup_{i\in I} U_i$$
    por lo que existen $i_1, \ldots, i_r \in I$ tales que $Y \subset (U_{i_1} \cup \cdots \cup U_{i_r})$, luego
    $$Y = Y \cap Y \subset (U_{i_1} \cup \cdots \cup U_{i_r}) \cap Y = \underbrace{(U_{i_1}\cap Y)}_{V_{i_1}} \cup \cdots \cup \underbrace{(U_{i_r}\cap Y)}_{V_{i_r}}$$
    por lo que $Y = V_{i_1} \cup \cdots \cup V_{i_r}$, lo que demuestra que $\{V_{i_1}, \ldots, V_{i_r}\}$ es un subrrecubrimiento finito de $Y$, por lo que $(Y,\mc T_X|_Y)$ es compacto. 
\end{itemize}
\end{dem}
\begin{prop}[Todo cerrado de un compacto es compacto]
Si $X$ es compacto y $Y \subset X$ es cerrado $\Rightarrow$ $Y$ es compacto. 
\end{prop}
\begin{dem}
Sean $\{U_i\}_{i\in I}$ abiertos de $X$ con $\displaystyle Y \subset \bigcup_{i\in I} U_i$, como $Y$ es cerrado, $X-Y$ es abierto, de forma que 
$$X = Y \cup (X-Y) \subset \bigcup_{i\in I} U_i \cup (X-Y) \subset X \Rightarrow X = \bigcup_{i\in I} U_i \cup (X-Y)$$
es un recubrimiento de $X$ por abiertos de $X$, pero por ser $X$ compacto, existe $\{i_1, \ldots, i_r \} \subset I$ tal que $X = U_{i_1} \cup \cdots \cup U_{i_r} \cup (X-Y)$ (puede que $X-Y$ aparezca o no en él; en cualquier caso, al añadirlo se mantiene la igualdad). Finalmente, dado cualquier $y \in Y$, se tiene $y \in X$ e $y \nin X-Y$, por lo que necesariamente $y \in U_{i_1} \cup \cdots \cup U_{i_r}$. Por tanto $Y \subset U_{i_1} \cup \cdots \cup U_{i_r}$ y se concluye que $Y$ es compacto. 
\end{dem}
\section{Compacidad y espacios de Hausdorff ($T_2$)}
\begin{prop}
Sea $X$ espacio topol\'ogico $T_2$ e $Y\subset X$ subespacio compacto. Entoces se verifica que para todo $x\nin Y$ existen abiertos $U_x, U_Y$ de $X$ tal que 
$$x\in U_x, \quad Y \subset U_Y \quad \text{con} \quad U_x\cap U_Y =\emptyset$$
\end{prop}
\begin{obs}
    Esta propiedad es análoga a la definición de espacio regular, pero separando un punto de un compacto en lugar de un cerrado arbitrario.
\end{obs}
\begin{dem}
Sea $x\nin Y$, dado $y\in Y$, claramente $x\neq y$ por ser $X$ separado, entonces existen abiertos 
$$x\in U_x^y, \quad y \in U_y \quad \text{con} \quad U_x^y \cap U_y =\emptyset$$
(nótese que el abierto $U_x^y$ depende tanto del punto $x$ como del punto $y\in Y$ considerado). Al variar $y$ en todo $Y$, la familia $\{U_y\}_{y \in Y}$ es una familia de abiertos de $X$ que recubre a $Y$, puesto que cada $y \in Y$ pertenece a su correspondiente $U_y$: 
$$Y\subset \bigcup_{y\in Y} U_y$$
y como $Y$ es compacto, entonces existen $y_1, \ldots, y_r \in Y$ tales que 
$$Y \subset U_{y_1}\cup \cdots \cup U_{y_r} \eqqcolon U_Y$$
Definimos
$U_x \coloneqq U_x^{y_1} \cap \cdots \cap U_x^{y_r}$ que es abierto por ser intersecci\'on finita. 
Comprobamos que cumplen todas las condiciones exigidas:
\begin{itemize}
    \item \textbf{Son abiertos de $X$:} $U_Y$ es abierto por ser unión de abiertos, y $U_x$ es abierto por ser \textbf{intersección finita} de abiertos.
    \item \textbf{Contienen a $x$ y a $Y$:} Por construcción, $Y \subseteq U_Y$. Además, como $x \in U_x^{y_j}$ para todo $j \in \{1, \dots, r\}$, se tiene $x \in \displaystyle \bigcap_{j=1}^r U_x^{y_j} = U_x$.
    \item \textbf{Son disjuntos:} $U_x\cap U_Y = \emptyset$. Sea $z\in U_x \cap U_Y$
    $$\Rightarrow \begin{cases}
        z\in U_x \Rightarrow z \in U_x^{y_j} \forall j=1,\ldots, r \\ z\in U_Y \Rightarrow z \in U_{y_{j_o}} \text{ para alg\'un } j_o \in \{1, \ldots, r\}
    \end{cases}$$
    luego $z\in U_x^{y_{j_o}}\cap U_{y_{j_o}} = \emptyset$, lo cual es una contradicci\'on al suponer que $z \in U_x \cap U_y$. 
\end{itemize}
\end{dem}
\begin{cor}[Todo compacto en un espacio $T_2$ es cerrado]
Si $X$ es $T_2$ e $Y\subset X$ es compacto $\Rightarrow Y$ es cerrado. En particular, si $X$ es $T_2$ y compacto, un subespacio $Y\subset X$ es compacto $\iff$ $Y$ es cerrado. 
\end{cor}
\begin{dem}
    $Y$ cerrado $\iff \, X-Y$ es abierto $\iff$ $\forall x\in X-Y$ existe un abierto $U_x$ con $x\in U_x \subset X-Y$ (entorno de todos sus puntos). 
    
    Sea $x\nin Y$, por hip\'otesis $Y$ es compacto, luego existen abiertos $U_x, UYy$ tales que $x\in U_x, Y \subset U_Y$ con $U_x\cap U_Y =\emptyset$ luego
    $$U_x\cap Y \subset U_x\cap U_Y = \emptyset \Rightarrow U_x\cap Y = \emptyset \Rightarrow U_x \subset X-Y$$ Así, todo punto $x \in X - Y$ posee un entorno abierto $U_x$ contenido en $X - Y$ (es decir, $X - Y = \bigcup_{x \in X - Y} U_x$), de modo que $X - Y$ es abierto y, por tanto, $Y$ es cerrado en $X$.

    Para la segunda afirmación, supongamos que $X$ es $T_2$ y compacto:
    \begin{itemize}
        \item[$\Rightarrow)$] Si $Y \subseteq X$ es compacto, entonces $Y$ es cerrado por lo que acabamos de demostrar (usando que $X$ es $T_2$).
        \item[$\Leftarrow)$] Si $Y \subseteq X$ es cerrado, entonces $Y$ es compacto por la Proposición 2 (usando que $X$ es compacto). \qedhere
    \end{itemize}
\end{dem}
\section{Aplicaciones continuas}
\begin{prop}[La imagen continua de un compacto es compacta]
    Si $f: X \longrightarrow Y$ es continua y $X$ es compacto entonces $f(X) = \Im \, f$ es compacto. 
\end{prop}
\begin{dem}
Veamos que $f(X)$ es un subespacio compacto de $Y$. Sea $\{V_i\}_{i \in I}$ una familia de abiertos de $Y$ tal que
\[
    f(X) \subseteq \bigcup_{i \in I} V_i.
\]
Tomando antiimágenes por $f$ y usando sus propiedades elementales
\[
    X = f^{-1}(f(X)) \subseteq f^{-1}\!\left(\bigcup_{i \in I} V_i\right) = \bigcup_{i \in I} f^{-1}(V_i) \subseteq X,
\]
y como $f^{-1}(V_i)\subset X$, resulta la igualdad $\displaystyle X = \bigcup_{i \in I} f^{-1}(V_i)$.

Como $f$ es \textbf{continua} y cada $V_i$ es abierto en $Y$, cada conjunto $f^{-1}(V_i)$ es abierto en $X$. Por tanto, $\{f^{-1}(V_i)\}_{i \in I}$ es un recubrimiento abierto de $X$.

Dado que $X$ es compacto, existen índices $\{i_1, \dots, i_r\} \subset I$ tales que:
\[
    X = f^{-1}(V_{i_1}) \cup \dots \cup f^{-1}(V_{i_r}) = \bigcup_{j=1}^r f^{-1}(V_{i_j}).
\]
Aplicando ahora $f$ a ambos miembros y recordando que $f(f^{-1}(B)) \subseteq B$ para todo subconjunto $B \subseteq Y$:
\[
    f(X) = f\!\left(\bigcup_{j=1}^r f^{-1}(V_{i_j})\right) = f\!\left(f^{-1}\!\left(\bigcup_{j=1}^r V_{i_j}\right)\right) \subseteq \bigcup_{j=1}^r V_{i_j}.
\]
Hemos encontrado así una subfamilia finita $\{V_{i_1}, \dots, V_{i_r}\}$ que recubre a $f(X)$, por lo que $f(X)$ es compacto.
\end{dem}
\begin{ejer}
$\bb R^n$ no es compacto $\forall n$. 
\end{ejer}\begin{dem}
Consideramos la proyecci\'on can\'onica $\pi_1: \bb R^n \longrightarrow \bb R$ que es continua (por la definici\'on de la topolog\'ia producto $\mc T_{\pi} = \mc T_d$ usual) y epiyectiva, luego $\pi_1(\bb R^n) = \bb R$.

Por la proposici\'on anterior la imagen continua de un compacto es compacta. Si $\bb R^n$ fuera compacto, su imagen por $\pi_1$, $\pi_1(\bb R^n) = \bb R$ debe ser compacto, pero sabemos que $\bb R$ no es compacto. Por tanto $\bb R^n$ no es compacto.  
\end{dem}